\section*{Chapitre \ref{ch:b3:organometallic-bonding} --- Chimie organométallique~: liaisons et ligands}

\begin{solution}{exo:b3:organometallic-bonding:1}
$\ce{Fe(CO)5}$~: $8 + 10 = 18$. $\ce{Mn2(CO)10}$~: $7 + 10 + 1 = 18$ par Mn. $\ce{CpMo(CO)3H}$~: $6 + 5 + 6 + 1 =
18$. $\ce{[PtCl3(C2H4)]-}$~: $10 + 3 + 2 + 1 = 16$. $\ce{Cr($\eta^6$-C6H6)(CO)3}$~: $6 + 6 + 6 = 18$. $\ce{Cp2ZrCl2}$~:
$4 + 10 + 2 = 16$.
\end{solution}

\begin{solution}{exo:b3:organometallic-bonding:2}
Fe(0) $d^8$~; Mn(0) $d^7$~; Mo(II) $d^4$~; Pt(II) $d^8$~; Cr(0) $d^6$~; Zr(IV) $d^0$.
\end{solution}

\begin{solution}{exo:b3:organometallic-bonding:3}
$\ce{[Mn(CO)6]+} > \ce{Cr(CO)6} > \ce{[V(CO)6]-}$~: plus le métal est riche en électrons, plus il
rétrocède dans $\pi^*$ et plus les liaisons C--O sont faibles.
\end{solution}

\begin{solution}{exo:b3:organometallic-bonding:4}
L'éthane ($\ce{Mn(CO)5}$ et $\ce{CpFe(CO)2}$ sont \omterm{def:b3:organometallic-bonding:isolobal}{isolobaux} à \ce{CH3})~; de nouveau l'éthane~; le tétraédrane
$\ce{(CH)4}$ dont trois CH sont remplacés par $\ce{Co(CO)3}$.
\end{solution}

\begin{solution}{exo:b3:organometallic-bonding:5}
$fac$ ($C_{3v}$)~: deux (A$_1$ + E). $mer$ ($C_{2v}$)~: trois (2A$_1$ + B$_1$).
\end{solution}

\begin{solution}{exo:b3:organometallic-bonding:6}
La phosphine encombrée surcharge le métal~: la dissociation soulage la
contrainte, si bien qu'elle est plus rapide, et les intermédiaires de
faible coordinence sont favorisés.
\end{solution}

\begin{solution}{exo:b3:organometallic-bonding:7}
Le complexe du chrome est un \omterm{def:b3:organometallic-bonding:carbene}{carbène de Fischer} (substituant OMe, Cr(0),
ligands CO)~: une amine attaque le carbone carbénique électrophile et
remplace OMe. L'alkylidène du tantale est un \omterm{def:b3:organometallic-bonding:carbene}{carbène de Schrock}~: son
carbone nucléophile attaque le carbonyle d'une cétone, ce qui donne un
alcène et un motif Ta=O, comme dans une réaction de Wittig.
\end{solution}

\begin{solution}{exo:b3:organometallic-bonding:8}
Décalé~: pas de recouvrement $\delta$, les deux électrons $\delta$ sont non liants~: indice de liaison 3. $\ce{[Re2Cl8]^{4-}}$~:
$\sigma^2\pi^4\delta^2\delta^{\ast2}$, indice de liaison 3.
\end{solution}

\begin{solution}{exo:b3:organometallic-bonding:9}
Une distance M$\cdots$H courte (diffraction, de préférence de neutrons)~;
un nombre d'onde d'élongation C--H bas~; un signal RMN à champ fort avec un
couplage C--H à une liaison réduit.
\end{solution}

\begin{solution}{exo:b3:organometallic-bonding:10}
16~: les complexes plans carrés $d^8$ ($\ce{[PtCl4]^{2-}}$, catalyseur de Wilkinson), dont l'orbitale du dix-huitième électron
est trop haute~; les \omterm{def:b3:organometallic-bonding:metallocene}{métallocènes} des métaux de début de série comme $\ce{Cp2ZrCl2}$, trop encombrés pour
accueillir d'autres ligands. 17~: $\ce{V(CO)6}$, qui ne peut pas se dimériser pour des raisons stériques. 19~: le cobaltocène,
dont l'électron supplémentaire occupe une orbitale antiliante (et se perd
facilement).
\end{solution}

\begin{solution}{exo:b3:organometallic-bonding:11}
CO est un fort accepteur $\pi$~: il abaisse le niveau $t_{2g}$ et rend $\Delta_{\mathrm o}$ grand, bien au-dessus de
l'énergie d'appariement~: $t_{2g}^6$, diamagnétique. Les transitions $d$--$d$, à $\Delta_{\mathrm o}$, tombent dans l'ultraviolet,
et aucune lumière visible n'est absorbée.
\end{solution}

\begin{solution}{exo:b3:organometallic-bonding:12}
La donation depuis $\pi$ et la rétrodonation dans $\pi^*$ abaissent toutes deux l'indice de liaison C--C. Dans le sel
de Zeise (Pt(II), rétrodonation modérée), la liaison s'allonge un peu~; un
métal à bas degré d'oxydation, riche en électrons, rétrocède fortement et
le complexe tend vers la limite du \omterm{def:b3:organometallic-bonding:dcd}{métallacyclopropane}, avec une liaison
C--C proche d'une liaison simple.
\end{solution}

\begin{solution}{pb:b3:organometallic-bonding:1}
\textbf{1.} $\ce{Fe^{2+}}$ ($d^6$) $+ 2 \times 6 = 18$.
\textbf{2.} $8 + 2 \times 5 = 18$.
\textbf{3.} $9 + 10 = 19$.
\textbf{4.} $10 + 10 = 20$.
\textbf{5.} $18 - 1 = 17$.
\textbf{6.} Fe(II) $d^6$~; Co(II) $d^7$~; Ni(II) $d^8$~; Fe(III) $d^5$.
\textbf{7.} $a_{1g}$ ($d_{z^2}$), $e_{1g}$ ($d_{xz}$, $d_{yz}$), $e_{2g}$ ($d_{xy}$, $d_{x^2-y^2}$).
\textbf{8.} La paire $e_{1g}$ recouvre fortement la combinaison $e_{1g}$ pleine des cycles~: des orbitales
liantes surtout centrées sur les cycles, et des antiliantes $e_{1g}^\ast$ surtout centrées sur le métal, hautes en
énergie.
\textbf{9.} $e_{2g} \approx a_{1g} < e_{1g}^\ast$.
\textbf{10.} $e_{2g}^4a_{1g}^2$~: aucun électron non apparié.
\textbf{11.} Un électron dans $e_{1g}^\ast$~: un électron non apparié, $1.73\,\mu_B$.
\textbf{12.} Deux électrons dans les deux orbitales $e_{1g}^\ast$, parallèles~: deux électrons non appariés, $\sqrt{8} =
2.83\,\mu_B$.
\textbf{13.} $e_{2g}^3a_{1g}^2$~: un électron non apparié. Un trou $e_{2g}^3$ donne un état orbitalement dégénéré
(E), avec une \omterm{def:b3:complex-spectra-magnetism:moment}{contribution orbitale}.
\textbf{14.} Son dix-neuvième électron occupe une orbitale antiliante et
s'arrache facilement, ce qui donne le cation cobaltocénium à 18 électrons.
\textbf{15.} L'électron arraché provient d'une orbitale presque non
liante~: la structure change à peine (faible \omterm{def:b3:complex-mechanisms:reorganisation}{énergie de réorganisation},
\cref{ch:b3:complex-mechanisms}).
\textbf{16.} Le couple est rapide et réversible dans la plupart des
solvants, le composé est soluble et stable, et son potentiel dépend peu du
solvant car la charge est répartie sur un gros ion.
\textbf{17.} Avec 20 électrons, dont deux antiliants, il réagit en les
perdant~: il fixe des ligands avec glissement ou perte d'un cycle, ou bien
il est oxydé.
\textbf{18.} $\ce{CpFe(CO)2}$~: $8 + 5 + 4 = 17$ électrons, un de moins que 18, une orbitale frontière~: isolobal
à \ce{CH3}. Le dimère $\ce{[CpFe(CO)2]2}$ est un «~éthane~», avec une liaison Fe--Fe (en pratique, deux ligands CO
la pontent).
\textbf{19.} $\ce{CpFe(CO)2-Mn(CO)5}$, avec une liaison Fe--Mn.
\textbf{20.} $\ce{PPh3}$ est meilleur donneur et moins bon accepteur $\pi$ que CO~: le métal rétrocède
davantage aux CO restants, dont les élongations descendent.
\textbf{21.} Deux (A$_1$ + E).
\textbf{22.} Le glissement du cycle vers $\eta^3$ libère l'équivalent de deux électrons de coordination~: un ligand
entrant peut se fixer de façon associative sans que le métal dépasse 18
électrons.
\textbf{23.} \textbf{$\mu_{\text{spin seul}}(\text{nickelocène}) = \sqrt{2 \times 4} = 2.83\,\mu_B$.}
\end{solution}
